% RESEARCH_DRAFT_SHA256: f9ecfe0eff9657ed422cbf9f70166b7968b3347769dc432e94a3d5de64a0aabc
% Layout and source-reference wrappers expanded for the web renderer.
\documentclass[11pt]{article}
\usepackage[margin=25mm]{geometry}
\usepackage{fontspec,unicode-math}
\setmainfont{Latin Modern Roman}
\setmathfont{Latin Modern Math}
\usepackage{amsmath,amssymb,mathtools,xurl,hyperref}
\hypersetup{colorlinks=true,urlcolor=blue,linkcolor=blue}
\setcounter{secnumdepth}{0}
\setlength{\parindent}{0pt}
\setlength{\parskip}{5pt}
\setlength{\emergencystretch}{3em}
\newcommand{\sref}[2]{\hyperlink{source:#1:#2}{[S#2]}}
\newcommand{\cataloguescope}[1]{\paragraph{Source record.}#1}
\newcommand{\floor}[1]{\left\lfloor#1\right\rfloor}
\begin{document}
\section*{Research-attempt provenance}
This standalone research attempt retains the complete original definition body below,
including its historical source and preparation notes. Those notes describe the
earlier catalogue preparation; the new attempt and its verification are identified
separately after that body. No original ID, catalogue status, or source statement is changed.

% UnsolvedMath ID 20001173; source catalogue code AIM-COMPUTATION-0011
\section{Finite-switch implementation of chattering controls}\label{20001173}
{\small UnsolvedMath ID 20001173 · Computer Science}
\subsection{Definitions and mathematical statement}

\paragraph{Shared notation.}Standard number systems are $\mathbb N,\mathbb Z,\mathbb Q,\mathbb R,\mathbb C$. Specialized notation and index ranges follow the cited source and the definitions below.


A feedback control is derived from state information, often through a value function solving an HJB equation. A chattering control switches infinitely often over a finite time interval. A finite-switch approximation replaces it by finitely many control intervals. Its performance loss compares the resulting cost with the optimal cost for the same system; discretization error additionally depends on the numerical approximation of the value function.
\paragraph{Record scope and qualifications.}
The catalogue title is retained as metadata. Named certificates, contracts, bounds or frameworks in that title are not asserted as results of the original source extract.
The excerpt contains several practical questions but no specified dynamics, switching budget, norm, or rate. Its qualitative finite-switch question is retained without a universal performance bound. The supplied catalogue marks this record partially solved; its present research status was not independently checked.
The source text below is rendered literally where necessary; damaged formulas and composite extracts are not silently repaired into a new mathematical claim.
\paragraph{Statement recorded in the supplied source.}
A.8 Oberman, Adam

I am interested in finite difference methods for high dimensional control and stochastic control problems. The types of general questions I am interested in are in 2 groups: - we know in theory that we can solve the PDE for the value function. How can we do this in practice for medium dimensional (3-4-5) problems? - Given an approximate numerical solution of the PDE for 9

the value function, it isn't trivial to implement a near-optimal control strategy. How do we best do this and can we understand how far from optimal these approximations are? I am curious about how much of a gap there is between theory and practice for nonlinear control problems. Given a control problem, we can write down the PDE for the value function, and we know that this PDE has unique viscosity solutions. From solutions of the PDE, we can then obtain the control. What I would like to know is the following: Suppose we have an approximate solution of the PDE, and we use this approximation to derive a control. How far from optimal will the resulting solution be? Also: suppose the optimal solution swtiches strategy infinitely many times. (This happens even in well-known examples, like stabilizing the linear/nonliner pendulum). Even supposing we had the exact solution, how far off are we from optimal if we approximate the control by one which only switches finitely many times, (which is all that we can do in practice). \href{https://aimath.org/WWN/optimalcontrol/optimalcontrol.pdf}{[S2]}
\paragraph{Source status.}
The supplied catalogue labels this source record partially solved; that metadata was checked against the supplied snapshot on 2026-10-11. The label is preserved for provenance and is not an independent verification of present open status or mathematical correctness.
\subsection{Short English statement}
The source asks how approximate value functions and finitely many switches affect control performance.
\subsection{Sources}
\begin{itemize}
\item[S1] Ulam AI and the UnsolvedMath Contributors, UnsolvedMath: A Curated Collection of Open Mathematics Problems, record 20001173 (AIM-COMPUTATION-0011). Catalogue curation and metadata: CC BY 4.0. The source extract is retained; definitions, English summaries and qualifications are editorial additions.
\url{https://huggingface.co/datasets/ulamai/UnsolvedMath}.
\item[S2] AIM Workshop Problem Lists, source linked by catalogue record AIM-COMPUTATION-0011.
\url{https://aimath.org/WWN/optimalcontrol/optimalcontrol.pdf}.
\end{itemize}
\paragraph{Preparation and review.}
The original source and UnsolvedMath contributors retain attribution for the source extract. Codex (GPT-6; exact serving revision not exposed) assisted with definitions, English summaries, scope notes and formatting. No new human mathematical authorship or independently verified solution is asserted. Permanent source ID 20001173 and code AIM-COMPUTATION-0011 checked against the complete supplied registry and current repository. Distinct upstream records, including repeated questions and broad or damaged extracts, are retained for complete coverage of the source CS category. A repeated record is not a new independent conjecture. Complete mathematical-equivalence and current-literature review has not been performed. Preparation and final review follow contribution prompts C and D, with source gaps disclosed. Structural and rendering checks do not establish a new theorem or open status.
\label{end:20001173}

\section{Literature and context check}
The participant contribution in the primary AIM optimal-control workshop
 PDF was consulted. Its broad scope is retained above. The following explicitly
 chosen finite or smooth model tests one route suggested by the agenda; it is
 not a reconstruction of a missing theorem. The elementary arguments are not
 claimed novel, and no current-resolution audit of the whole field is asserted.
\section{Research attempt: A finite-switch truncation bound depends on the switching tail}
Choose dynamics $\dot x=f(x,u)$ on a fixed horizon $[0,T]$, assuming the
original and truncated-control trajectories both exist throughout that horizon, with
$\|f\|\le M$ and state Lipschitz constant $L$, uniformly over admissible
actions. Let a given chattering control have switching times
$0<t_1<t_2<\cdots\uparrow T$. Retain it until $t_N$ and hold one admissible
action afterward. Let $x,\widetilde x$ be the original and truncated trajectories.
They agree before $t_N$. Afterwards
\[
 \|x(t)-\widetilde x(t)\|\le
 \begin{cases}2M(e^{L(t-t_N)}-1)/L,&L>0,\\2M(t-t_N),&L=0.\end{cases}
\]
The proof subtracts the dynamics, bounds the action-change term by $2M$ and
integrates $e'\le Le+2M$. In particular the uniform error is at most
$2Me^{LT}(T-t_N)$, using $e^{Ls}-1\le Lse^{LT}$.

Suppose $|c(x,u)|\le B$, $c$ is $K$-Lipschitz in state uniformly in $u$,
and terminal cost is $K_g$-Lipschitz. The running-cost difference is bounded
by $K\|x-\widetilde x\|$ from state change plus $2B$ only on the altered
tail. Hence
\[
 |J(\widetilde u)-J(u)|
 \le[2M e^{LT}(KT+K_g)+2B](T-t_N).
\]
If the given chatter is optimal, this is a suboptimality bound for the finite
switch control. If it is not known optimal, the same formula only compares
these two controls; its optimality gap must be added separately.

The attempted universal rate in $N$ cannot be read from this argument.
For any decreasing positive sequence $a_N\to0$ with $a_0=T$, times
$t_N=T-a_N$ give a valid accumulating schedule (strictly decrease $a_N$ to
ensure distinct switches). Thus the proved bound can decay geometrically,
polynomially, or more slowly depending on the actual tail. This is not a
lower bound for every best $N$-switch approximation; it precisely identifies
the extra assumption needed by the proposed truncation proof. Smooth dynamics
alone do not specify those switching times.

The other part of the source, deriving feedback from a numerical value
function, introduces derivative and implementation errors not estimated here.
\section{Verification and adversarial audit}
The comparison has identical controls before truncation and includes
both direct action-cost changes and terminal cost. The $L=0$ case was derived
without division by zero. The schedule argument is only about what this bound
can certify, not an unproved impossibility result for optimized approximations.
\paragraph{Audit corrections.} A comparison to a given chatter was called suboptimality only conditionally on that chatter being optimal; no universal rate or pendulum result was invented.
\paragraph{External computation.} None; symbolic arguments and boundary checks only.
\section{Scope and remaining gaps}
Actual dynamics, switch accumulation law, costs and numerical scheme remain
unspecified. The first gap for a rate in switching budget is quantitative
control of the tail or a separate optimal finite-switch construction, plus
the missing value-to-feedback error analysis.
\paragraph{Integration review.} Codex (GPT-6; exact serving revision not exposed) reviewed the written arguments and expanded layout and source-reference wrappers. Repository metadata and the production web renderer were checked using Python 3.13 and JavaScript with MathJax 3.2.2. This model-assisted review is not independent human proof verification; no PDF compilation or formal proof checking is asserted.

\section{Completion Estimate}
\noindent\textbf{COMPLETION: 0\%}
The calculations below the preserved source statement investigate a concrete route,
but establish no substantiated new progress toward the entire original target.
Known special cases and failed sufficient conditions are not credited as a solution.
\section{References}
The original source references and catalogue attribution are preserved in the
Sources section above. The following URLs were actually consulted for this attempt;
retention of an additional original reference does not claim fresh consultation of it.
\begin{itemize}
\item \url{https://aimath.org/WWN/optimalcontrol/optimalcontrol.pdf} (consulted 2026-10-11). 
\end{itemize}
\section{Model and contributor disclosure}
OpenAI gpt-6-astra, requested reasoning configuration ultra, generated this research
attempt on 2026-10-11. The exact serving revision is not exposed. Codex GPT-6
performed orchestration, repository-format review, and a model-assisted mathematical
audit of the written arguments; its exact serving revision
is also not exposed. The mathematical route, boundary checks and final adversarial
self-audit followed contribution prompts A then B in
\texttt{docs/contribution\_prompts.md}, after reading \texttt{CONTRIBUTING.md}.
Python 3.14.2 assembled the standalone documents; any mathematical computation is
specified in Verification. No theorem prover was used. No human mathematical
authorship, substantial human proof correction, or independent human proof checking
is asserted. Model-assisted review, source retrieval, and structural checks do not
establish correctness or novelty. Full private reasoning transcripts are unavailable
as reproducibility artifacts; the complete arguments, inputs and outcomes actually
used are stated here or in the companion batch audit.
\section{Final status}
Unresolved as a new research contribution. This is a concrete route investigation,
not a claim of new complete resolution or a current-status certification. A known
answer reproduced for context does not become a new contribution. The catalogue's
source-status label is preserved; this attempt status does not assert that a
source-labelled solved question is open.
\end{document}
